\section{Development in the IDE}\label{sec:ide}

\subsection{Setting up the project}

The project was created as a Maven project with the Maven wrapper and opened in Visual Studio Code as a folder. The
Java extension reads \texttt{pom.xml}, imports the project with its dependencies, and compiles it in the background;
errors and warnings appear in the editor as the code is typed. Fig.~\ref{fig:ide-01-project} shows the project open,
its folders in the Explorer and the integrated terminal below the editor.

\screenshotwide{ide-01-project}
  {The project open in Visual Studio Code: the project tree in the Explorer, and the integrated terminal after a run of
   the Javadoc tool.}
  {VS Code with the project open: Explorer tree, Java Projects view, JDK 25}

The folder holds the Maven wrapper (\texttt{mvnw}, \texttt{mvnw.cmd} and \path{.mvn/}), \texttt{pom.xml}, the
IDE's settings in \path{.vscode/}, and the README. The code is under \path{src/main/java}, its resources (the
settings, the Flyway migration, and the window's FXML, stylesheet and fonts) under \path{src/main/resources}, and the
tests under \path{src/test/java}. The generated documentation is in \path{docs/javadoc/}, and this report in
\path{report/}.

\subsection{Running the program}

The project carries a launch configuration for the IDE, ``Task Manager'' in \path{.vscode/launch.json}: it starts
the main class \texttt{io.mcs.taskmanager.TaskManagerApplication} on JDK~25 with
\texttt{--enable-native-access=ALL-UNNAMED}, from the Run and Debug view or with F5, and shows the program's output in
the Debug Console. The program also starts from the IDE's integrated terminal with the Maven wrapper,
\texttt{.\textbackslash{}mvnw.cmd spring-boot:run} in PowerShell or \texttt{./mvnw spring-boot:run} in Git Bash.
Fig.~\ref{fig:ide-02-run} shows it started that way: its window, on the empty list, open over the editor, and
Spring's start-up line in the terminal behind it.

\screenshotwide{ide-02-run}
  {The program started from the IDE's integrated terminal with \texttt{.\textbackslash{}mvnw.cmd spring-boot:run}: its window, on
   the empty list, over the editor.}
  {The app started from the terminal}

\subsection{Debugging}

Because the launch configuration is of type \texttt{java}, the IDE starts the program under the Java debugger: a
breakpoint set in the editor, for example on the first line of \texttt{TaskService.create}, stops the program there
when a task is added, and the Variables and Call Stack views then show its state, from which it can be stepped
through line by line.
